% !TEX program = lualatex
% Kompilacja: lualatex main.tex  (zalecane)
%         lub: pdflatex main.tex (wymaga pakietu metropolis i fontów)
% Jeśli nie masz metropolis: zainstaluj texlive-latex-extra / mactex
% Szybki fallback: zamień 'metropolis' na 'Madrid' w \usetheme

\documentclass[aspectratio=169,11pt]{beamer}

% ---------- język i fonty ----------
\usepackage[T1]{fontenc}
\usepackage[utf8]{inputenc}
\usepackage[polish]{babel}
\usepackage{lmodern}
\usepackage{microtype}

% ---------- motyw ----------
\usetheme{metropolis}
\metroset{
  progressbar=frametitle,
  numbering=fraction,
  block=fill,
  sectionpage=progressbar,
  subsectionpage=progressbar
}

% ---------- kolory (akcent) ----------
\definecolor{mPrimary}{HTML}{1F4E79}   % granat akademicki
\definecolor{mAccent}{HTML}{E07A1F}    % pomarańcz – akcent
\setbeamercolor{progress bar}{fg=mAccent}
\setbeamercolor{title separator}{fg=mAccent}
\setbeamercolor{frametitle}{bg=mPrimary}
\setbeamercolor{alerted text}{fg=mAccent}

% ---------- pakiety pomocnicze ----------
\usepackage{booktabs}
\usepackage{array}
\usepackage{tabularx}
\usepackage{graphicx}
\usepackage{tikz}
\usepackage{pgfgantt}
\usepackage{hyperref}
\usepackage{appendixnumberbeamer}

\hypersetup{
  colorlinks=true,
  linkcolor=mPrimary,
  urlcolor=mAccent,
  citecolor=mPrimary
}

% ---------- metadane ----------
\title{Prywatność na pierwszym miejscu}
\subtitle{Aplikacja mobilna z lokalnymi modelami językowymi i rozpoznawaniem mowy w~React~Native}
\author{Nikodem Nowak}
\institute{Wydział ... \\ Kierunek ...}
\date{Proseminarium \textbullet{} kwiecień 2026}

% ==============================================================
\begin{document}

% ---------- tytuł ----------
\begin{frame}[plain]
  \titlepage
\end{frame}

% ---------- agenda ----------
\begin{frame}{Plan prezentacji}
  \tableofcontents
\end{frame}

% ==============================================================
\section{Motywacja}

\begin{frame}{Dlaczego ten temat?}
  \begin{itemize}
    \item LLM-y weszły do codziennego użycia -- ChatGPT, Gemini, Copilot.
    \item Każde zapytanie = dane użytkownika opuszczają urządzenie.
    \item Wrażliwe konteksty: medycyna, prawo, finanse, dziennikarstwo.
    \item Regulacje: \alert{RODO}, \alert{AI~Act} (obowiązuje od 2024/2025).
    \item Trend ,,local-first AI'': Apple Intelligence, Gemini Nano, Llama~3.2, Phi-3.
  \end{itemize}
  \vfill
  \begin{alertblock}{Obserwacja}
    Modele stały się na tyle małe, a telefony na tyle mocne, że pełna prywatność \emph{bez~kompromisu na funkcjonalności} jest dziś realna.
  \end{alertblock}
\end{frame}

\begin{frame}{Problem badawczy}
  \begin{center}
    \Large
    Czy współczesne urządzenia mobilne są gotowe,\\
    by w~pełni \alert{lokalnie} obsłużyć asystenta głosowego z~LLM-em?\\[1em]
    \normalsize
    I~które rozwiązanie runtime'owe sprawdza się najlepiej?
  \end{center}
\end{frame}

% ==============================================================
\section{Teza i cele}

\begin{frame}{Hipotezy}
  \begin{block}{H1 -- produktowa}
    Możliwe jest zbudowanie użytecznej aplikacji mobilnej z~lokalnym STT i~LLM, działającej w~pełni offline, na~urządzeniach konsumenckich z~ostatnich 2--3 lat.
  \end{block}
  \vspace{0.5em}
  \begin{block}{H2 -- badawcza}
    Wybór runtime'u (\textit{ExecuTorch} vs \textit{MLC~LLM} vs \textit{llama.cpp}) istotnie wpływa na~latencję, zużycie pamięci i~energii przy~porównywalnej jakości wyjścia.
  \end{block}
\end{frame}

\begin{frame}{Cele pracy}
  \begin{enumerate}
    \item \textbf{Zaprojektować} architekturę aplikacji RN działającej w~pełni on-device.
    \item \textbf{Zaimplementować} działający produkt: głos~$\to$~transkrypcja~$\to$~LLM~$\to$~odpowiedź.
    \item \textbf{Zintegrować} co najmniej 2--3 runtime'y LLM z~jednolitym API.
    \item \textbf{Zmierzyć i~porównać} wydajność, zużycie zasobów i~jakość wyjścia.
    \item \textbf{Udokumentować} wnioski w~formie rekomendacji dla developerów RN.
  \end{enumerate}
\end{frame}

\begin{frame}{Zakres pracy}
  \begin{columns}[T,onlytextwidth]
    \column{0.48\textwidth}
    \textbf{W~zakresie}
    \begin{itemize}
      \item Android oraz iOS
      \item Modele LLM do~$\sim$4B parametrów
      \item Kwantyzacja 4-bit / 8-bit
      \item Whisper (\textit{base}, \textit{small})
      \item Ewaluacja ilościowa i~jakościowa
    \end{itemize}

    \column{0.48\textwidth}
    \textbf{Poza zakresem}
    \begin{itemize}
      \item Trening i~fine-tuning modeli
      \item Multimodalność (obraz, wideo)
      \item Rozwiązania chmurowe
      \item Wsparcie urządzeń starszych niż 2023
    \end{itemize}
  \end{columns}
\end{frame}

% ==============================================================
\section{Kontekst teoretyczny}

\begin{frame}{Jak LLM działa na telefonie?}
  \begin{itemize}
    \item \textbf{Kwantyzacja} -- zmniejszenie precyzji wag (FP16~$\to$~INT4/INT8).
          Model 3B z~6~GB spada do~$\sim$1.5~GB.
    \item \textbf{Formaty wag}: GGUF (llama.cpp), ExecuTorch \texttt{.pte}, MLC \texttt{.mlc}.
    \item \textbf{Akceleracja}: CPU z~SIMD, GPU (Metal/Vulkan), NPU (Core~ML, QNN, NNAPI).
    \item \textbf{Pamięć}: \texttt{mmap} wag -- system ładuje tylko potrzebne fragmenty.
    \item \textbf{STT}: Whisper skwantyzowany działa real-time na~CPU telefonu z~2023+.
  \end{itemize}
\end{frame}

\begin{frame}{Krajobraz rozwiązań -- RN / Expo}
  \small
  \begin{tabularx}{\textwidth}{@{}l l X@{}}
    \toprule
    \textbf{Biblioteka} & \textbf{Silnik} & \textbf{Charakterystyka} \\
    \midrule
    \texttt{react-native-executorch} & ExecuTorch (Meta) & Deklaratywne API, wsparcie LLM + Whisper + CV. \\
    \texttt{react-native-ai}          & MLC LLM (CMU)     & Kompatybilne z~Vercel AI SDK, wielorun\-time'owe. \\
    \texttt{llama.rn}                 & llama.cpp         & Niskopoziomowe, maksymalna kontrola, GGUF. \\
    \texttt{react-native-rag}         & ExecuTorch        & Lokalny RAG -- retrieval po własnych danych. \\
    \bottomrule
  \end{tabularx}
  \vspace{1em}
  \begin{itemize}
    \item Wszystkie trzy runtime'y są aktywnie rozwijane (2025--2026).
    \item Brak istniejących publikacji porównujących je na~jednym zestawie urządzeń $\Rightarrow$ \alert{luka badawcza}.
  \end{itemize}
\end{frame}

\begin{frame}{Architektura: chmura vs on-device}
  \centering
  \begin{tikzpicture}[
      node distance=1.2cm,
      box/.style={draw, rounded corners, minimum width=2.2cm, minimum height=0.9cm, align=center, font=\small},
      cloud/.style={draw, cloud, cloud puffs=12, minimum width=2cm, minimum height=1cm, font=\small},
      arr/.style={-latex, thick}
    ]
    % ---- chmura ----
    \node[box] (u1) {Użytkownik};
    \node[box, right=of u1] (app1) {Aplikacja};
    \node[cloud, right=of app1, fill=red!10] (cloud) {Chmura};
    \draw[arr, red!70!black] (u1) -- node[above,font=\scriptsize]{głos / tekst} (app1);
    \draw[arr, red!70!black] (app1) -- node[above,font=\scriptsize]{\textbf{dane}} (cloud);
    \node[above=0.1cm of cloud, font=\scriptsize\bfseries, red!70!black] {dane opuszczają urządzenie};

    % ---- on-device ----
    \node[box, below=1.8cm of u1] (u2) {Użytkownik};
    \node[box, right=of u2] (app2) {Aplikacja};
    \node[box, right=of app2, fill=green!10] (llm) {LLM\\ on-device};
    \draw[arr, green!50!black] (u2) -- (app2);
    \draw[arr, green!50!black] (app2) -- (llm);
    \draw[arr, green!50!black] (llm.south) .. controls +(0,-0.6) and +(0,-0.6) .. (app2.south);
    \node[below=0.1cm of llm, font=\scriptsize\bfseries, green!50!black] {zamknięta pętla};
  \end{tikzpicture}
\end{frame}

% ==============================================================
\section{Metodologia}

\begin{frame}{Stos technologiczny}
  \begin{columns}[T,onlytextwidth]
    \column{0.48\textwidth}
    \textbf{Frontend / aplikacja}
    \begin{itemize}
      \item React Native + Expo
      \item TypeScript (strict)
      \item Expo SDK (audio, FS, haptics)
    \end{itemize}

    \column{0.48\textwidth}
    \textbf{AI / modele}
    \begin{itemize}
      \item Whisper (\textit{base}, \textit{small})
      \item Llama~3.2 1B/3B, Qwen~3, Phi-3-mini
      \item \texttt{react-native-executorch}
      \item \texttt{react-native-ai} (MLC)
      \item \texttt{llama.rn}
    \end{itemize}
  \end{columns}
\end{frame}

\begin{frame}{Metodologia ewaluacji}
  \textbf{Metryki wydajności}
  \begin{itemize}
    \item TTFT -- \emph{time to first token}
    \item Przepustowość generacji [tok/s]
    \item RAM peak, rozmiar modelu na dysku
    \item Zużycie energii [\%~baterii / min], temperatura obudowy
  \end{itemize}
  \textbf{Metryki jakości}
  \begin{itemize}
    \item STT: WER (Word Error Rate) na~zbiorze Common Voice PL/EN
    \item LLM: subset MMLU-lite, MT-Bench-lite, własne prompty domenowe
  \end{itemize}
  \textbf{Urządzenia testowe (planowane)}
  \begin{itemize}
    \item Android: Pixel 8, średniopółkowiec (np.~Samsung~A55)
    \item iOS: iPhone 15~Pro, iPhone 13
  \end{itemize}
\end{frame}

\begin{frame}{Scenariusz użycia -- MVP}
  \centering
  \begin{tikzpicture}[
      node distance=0.9cm and 1.2cm,
      step/.style={draw, rounded corners, minimum width=2.4cm, minimum height=0.9cm, align=center, font=\small, fill=mPrimary!10},
      arr/.style={-latex, thick, mPrimary}
    ]
    \node[step] (s1) {Nagranie\\ głosu};
    \node[step, right=of s1] (s2) {Whisper\\ (STT)};
    \node[step, right=of s2] (s3) {LLM\\ lokalnie};
    \node[step, right=of s3] (s4) {Odpowiedź\\ + zapis};
    \draw[arr] (s1) -- (s2);
    \draw[arr] (s2) -- (s3);
    \draw[arr] (s3) -- (s4);
  \end{tikzpicture}
  \vspace{1.5em}
  \begin{itemize}
    \item Przykładowe zastosowanie: prywatny notatnik głosowy z~automatycznym podsumowywaniem i~wyciąganiem zadań.
    \item Opcjonalnie: RAG po notatkach użytkownika (\texttt{react-native-rag}).
  \end{itemize}
\end{frame}

% ==============================================================
\section{Harmonogram i ryzyka}

\begin{frame}{Harmonogram}
  \centering
  \begin{ganttchart}[
      hgrid, vgrid,
      x unit=0.95cm,
      y unit title=0.6cm,
      y unit chart=0.55cm,
      title/.style={fill=mPrimary!20, draw=mPrimary},
      title label font=\scriptsize\bfseries,
      bar/.style={fill=mAccent!70, draw=mAccent!90!black},
      bar label font=\scriptsize,
      milestone/.style={fill=mPrimary, draw=mPrimary},
      milestone label font=\scriptsize
    ]{1}{8}
    \gantttitlelist{"V","VI","VII","VIII","IX","X","XI","XII"}{1} \\
    \ganttbar{Research i~POC}{1}{2} \\
    \ganttbar{Implementacja MVP}{2}{4} \\
    \ganttbar{Integracja 3 runtime'ów}{4}{5} \\
    \ganttbar{Pomiary, ewaluacja}{5}{6} \\
    \ganttbar{Analiza wyników}{6}{7} \\
    \ganttbar{Redakcja pracy}{6}{8} \\
    \ganttmilestone{Obrona}{8}
  \end{ganttchart}
  \vspace{0.5em}
  \scriptsize{Miesiące odpowiadają 2026; obrona planowana na~grudzień.}
\end{frame}

\begin{frame}{Ryzyka i~plan łagodzenia}
  \begin{tabularx}{\textwidth}{@{}l X X@{}}
    \toprule
    \textbf{Ryzyko} & \textbf{Opis} & \textbf{Mitygacja} \\
    \midrule
    Breaking changes & Biblioteki młode, API się zmienia. & Pin wersji, fork na~stabilnej wersji. \\
    Limity pamięci iOS & iOS ubija aplikacje >~3~GB RAM. & Ograniczenie do~modeli $\leq$3B, testy OOM. \\
    Hałas pomiaru baterii & Różnice temperatury, stanu urządzenia. & Protokół: n~powtórzeń, stała temperatura startowa. \\
    Dostępność urządzeń & Dostęp do~konkretnych telefonów. & Plan B: emulatory + 2 realne urządzenia. \\
    \bottomrule
  \end{tabularx}
\end{frame}

% ==============================================================
\section{Podsumowanie}

\begin{frame}{Co mam zrobione / co dalej}
  \begin{itemize}
    \item[\textbullet] Research ekosystemu RN on-device AI -- \alert{zrobione}.
    \item[\textbullet] Wybór kierunku: apka (privacy-first asystent) + benchmark runtime'ów.
    \item[$\to$] \textbf{Następny krok:} proof-of-concept z~\texttt{react-native-executorch} + Whisper na~Pixelu 8.
    \item[$\to$] Kolejne proseminarium: demo pierwszej wersji aplikacji.
  \end{itemize}
\end{frame}

\begin{frame}[standout]
  Dziękuję.\\[0.5em]
  \large Pytania?
\end{frame}

% ==============================================================
\appendix

\begin{frame}[allowframebreaks]{Bibliografia -- wybór}
  \footnotesize
  \begin{thebibliography}{9}
    \bibitem{executorch}
      Meta AI, \emph{ExecuTorch: On-device AI inference framework}, 2024.
      \url{https://pytorch.org/executorch/}
    \bibitem{rnexecutorch}
      Software Mansion, \emph{react-native-executorch}, 2024--2026.
      \url{https://github.com/software-mansion/react-native-executorch}
    \bibitem{rnai}
      Callstack, \emph{react-native-ai}, 2025--2026.
      \url{https://github.com/callstackincubator/ai}
    \bibitem{mlcllm}
      MLC Team, \emph{MLC LLM: Universal LLM deployment engine}, CMU, 2023--2026.
    \bibitem{whisper}
      A.~Radford et~al., \emph{Robust Speech Recognition via Large-Scale Weak Supervision} (Whisper), OpenAI, 2022.
    \bibitem{llama3}
      Meta AI, \emph{The Llama~3 Herd of Models}, 2024.
    \bibitem{aiact}
      Parlament Europejski, \emph{Regulation (EU) 2024/1689 (AI~Act)}, 2024.
  \end{thebibliography}
\end{frame}

\end{document}
